% !PRINTED 173
\weberpage{173}
%<<CONT
n\"ahert sich $\varkappa_\nu$ mit unendlich wachsendem $\nu$ der Grenze Null. Daraus erh\"alt man f\"ur $q$ die Darstellung:

\weq{22}{q = \mathop{\mathrm{Lim}}\limits_{\nu=\infty} \sqrt[2^{\nu-1}]{\frac{\varkappa_\nu}{4\varkappa_\nu'}} = \mathrm{Lim} \sqrt[2^{\nu-1}]{\frac{\varkappa_\nu}{4}},}

die bei reellen Werten von $\varkappa$ zur Berechnung geeignet ist.

Aus jeder der Formeln (21) kann man eine Reihenentwickelung von $q$ nach den steigenden Potenzen von $\varkappa_\nu$ herleiten, deren Konvergenz mit wachsendem $\nu$ zunimmt. Wendet man insbesondere die Formel (20) an, so erh\"alt man eine Entwickelung, die von \gesperrt{Weierstrass} in den Monatsberichten der Berliner Akademie vom Jahre 1883 mitgeteilt ist, deren Konvergenzbereich die Ebene $\varkappa^2$ zweimal ausf\"ullt.

Die Koeffizienten dieser Entwickelung berechnet man am einfachsten aus:

\weq{23}{\sqrt{\varkappa_2} = \frac{\vartheta_{10}(0,4\omega)}{\vartheta_{00}(0,4\omega)} = \frac{2q+2q^9+2q^{25}+\cdots}{1+2q^4+2q^{16}+\cdots},}

und erh\"alt so nach der Methode der unbestimmten Koeffizienten:

\weq{24}{q = \frac{\sqrt{\varkappa_2}}{2} + 2\left(\frac{\sqrt{\varkappa_2}}{2}\right)^5 + 15\left(\frac{\sqrt{\varkappa_2}}{2}\right)^9 + 150\left(\frac{\sqrt{\varkappa_2}}{2}\right)^{13} + 1707\left(\frac{\sqrt{\varkappa_2}}{2}\right)^{17} + \cdots}

Ebenso l\"a\ss{}t sich nach (23) auch umgekehrt $\sqrt{\varkappa_2}$ in eine nach Potenzen von $q$ fortschreitende Reihe entwickeln:

\weq{25}{\tfrac{1}{2}\sqrt{\varkappa_2} = q - 2q^5 + 5q^9 - 10q^{13} + 18q^{17}\cdots,}

die man auch durch Umkehrung der Reihe (24) erh\"alt.

Damit ist die am Anfang dieses Paragraphen gestellte Frage vollst\"andig beantwortet.

Bez\"uglich der Anwendung der hier entwickelten Formeln zu numerischer Berechnung von $q$ sei noch bemerkt, da\ss{}, wenn $\varkappa^2$ n\"aher an 1 liegt, man ein besser konvergentes Verfahren erh\"alt, wenn man $\varkappa^2$ mit $\varkappa'^2$ vertauscht. Die Rechnung ergibt dann zun\"achst nicht $q$, sondern

\[
q' = e^{-\frac{\pi i}{\omega}},
\]

woraus man aber $q$ aus der Formel erh\"alt:

\[
\log q \log q' = \pi^2.
\]
